\section{Log-only remainder ledger}
\label{sec:logonly-remainder-ledger}

This is the remainder ledger used by the conditional theorem.  The entire
packet-polarized logarithmic source, including its finite-height dependence and
the additive carrier $e(-ne^d)$, belongs to $A_T^{\rm true}$.  Thus no scalar
approximation error and no frozen-profile difference is placed in the regular
remainder.  The remaining theorem-facing classes are measured with the
absolute-lag primitive criterion of
Section~\ref{sec:primitive-remainder-core}.

The exact right-edge reflection also contributes the Hermitian archimedean
self block $J_T^{\rm arch}=I_f+I_f^*$ of \eqref{eq:J-I-source-split}.
Lemma~\ref{lem:holomorphic-self-elimination} does not remove it. The following
identity keeps this block visible before any small-remainder hypothesis.

There is no direct-residual operator in this ledger.  The safe-line estimate
$\Delta H_{\rm dir}=O(L^{-2})$ does not survive the one-$\chi$ stationary
normalization uniformly in the relevant Dirichlet carrier.  The full source
$\mathcal C(H_{\log}+\Delta H_{\rm dir})=\mathcal C H_{\rm ar}$ belongs to
$A_T^{\rm true}$ by definition \eqref{eq:true-logonly-operator}.

\begin{lemma}[Log-only $f'$ class is relative-regular]
\label{lem:logonly-fprime-relative-regular}
Let $R_{f',T}^{\log}:=\mathfrak R_T[R_{f'}^{\log};1,1]$ be the Hermitian
right-edge contribution of \eqref{eq:logonly-fprime-source}. Then
\begin{equation}
\boxed{
\left\|\Gfr_T^{-1/2}R_{f',T}^{\log}\Gfr_T^{-1/2}\right\|_{S_2}^2=o(N_T).
}
\label{eq:logonly-fprime-relative-HS}
\end{equation}
\end{lemma}

\begin{proof}
Every term in \eqref{eq:logonly-fprime-source} contains at least one factor
$f'=O(T^{-1})$ on the retained height strip, while
$(f-P)^{-1}=O(L^{-1})$ on the fixed safe line.  Fixed stationary, shell,
curvature and Cauchy derivatives cost only a fixed power of $L$, so each
shell-local profile and its first two lag derivatives are $O(T^{-1}L^C)$.
The shell normalization and primitive criterion therefore give a normalized
jump $o(1)$ and total mixed-derivative Hilbert--Schmidt mass $o(N_T)$.
Moreover the boundary coefficients of
\eqref{eq:primitive-boundary-coefficients} satisfy
\[
\sup_j\|\mathcal B_{j,0}\|_{\op}\ll T^{-1}L^C=o(1),
\qquad
\sum_j\int_{I_j}(\|\mathcal B_{j,R}\|_{S_2}^2
+\|\mathcal B_{j,L}\|_{S_2}^2)
\ll T^{-2}L^CN_T=o(N_T).
\]
Thus the boundary-complete criterion applies without local mean deletion.
\end{proof}

For the final ledger write $R_T^{f'}:=R_{f',T}^{\log}$.

\begin{proposition}[Exact four-block contour decomposition]
\label{prop:exact-four-block-decomposition}
On the full packet space, define the archimedean self block and the exact
horizontal block by
\begin{align}
J_T^{\rm arch}
&:=\mathfrak R_T\bigl[f(\mathcal C+f'\zeta^2);1,1\bigr],
\label{eq:exact-arch-block}\\
\langle R_T^{\rm hor,H}v,w\rangle
&:=\frac1{2\pi i}\int_{\Gamma_{\rm hor}}
H(s)C_{\mathscr Z}^{\rm norm}(s)\mathcal B_{v,w}(s)\,ds,
\label{eq:exact-horizontal-H}
\end{align}
where $\Gamma_{\rm hor}$ is the union of the top and bottom sides with the
boundary orientation of $\partial\mathcal R_T^{\rm sym}$. Set
\begin{equation}
R_T^{\rm exact}:=R_{f',T}^{\log}+R_T^{\rm hor,H}.
\label{eq:exact-fourth-block}
\end{equation}
With $A_T^{\rm true}$ and $L_T^{\rm geom}$ as in
\eqref{eq:true-logonly-operator}--\eqref{eq:exact-geometric-leakage}, the
original normalized contour matrix satisfies
\begin{equation}
\boxed{
A_F=A_T^{\rm true}+L_T^{\rm geom}+J_T^{\rm arch}+R_T^{\rm exact}.
}
\label{eq:exact-four-block-decomposition}
\end{equation}
This equality is exact before stationary phase, shell estimates, compression,
or any asymptotic limit; it also holds after restricting all four terms to the
same retained subspace.
\end{proposition}

\begin{proof}
On the right edge, \eqref{eq:curvature-source} and
\eqref{eq:L1-logonly-split} give the algebraic identity
\[
H(s)C_{\mathscr Z}^{\rm norm}(s)
=\frac{q(s)^2}{L^2}\left[
\mathcal C(s)H_{\rm ar}(s)
+f(s)\bigl(\mathcal C(s)+f'(s)\zeta(s)^2\bigr)
+R_{f'}^{\log}(s)\right].
\]
Reflection \eqref{eq:right-left-matrix-adjoint} turns the sum of the right and
left integrals into the Hermitian right-edge operator $\mathfrak R_T$ applied
to these three sources.  The horizontal integral is exactly
\eqref{eq:exact-horizontal-H}. Moreover
$\mathfrak R_T[\mathcal C H_{\rm ar};1,1]
=A_T^{\rm true}+L_T^{\rm geom}$ by
\eqref{eq:exact-geometric-leakage}.  Finally, the closed contour integrals
with $H=f+L_1$ and with $L_1$ are equal by
\eqref{eq:same-form-identity}. This proves the four-block identity.
\end{proof}

\begin{hypothesis}[Geometric edge and regular-remainder bounds]
\label{hyp:edge-regular-bounds}
For the same packet data and retained space $\Pret$ as in
Hypothesis~\ref{hyp:one-chi-relative-defect}, assume there is an edge operator
$E_T^{\rm geom}$ such that
\begin{equation}
\boxed{
L_T^{\rm geom}+J_T^{\rm arch}
=E_T^{\rm geom}+R_T^{\rm edge+arch},
}
\label{eq:logonly-final-decomposition}
\end{equation}
with all operators restricted to $\Pret$. Here $A_T^{\rm true}$ is the exact
log-only operator \eqref{eq:true-logonly-operator},
\begin{equation}
\boxed{\rank E_T^{\rm geom}=o(N_T),}
\label{eq:logonly-edge-rank}
\end{equation}
and each exact component satisfies
\begin{equation}
\boxed{
\left\|\Gfr_T^{-1/2}R_T^\nu\Gfr_T^{-1/2}\right\|_{S_2}^2=o(N_T)
\quad(\nu\in\{\mathrm{edge+arch},f',\mathrm{hor}\}).
}
\label{eq:logonly-component-relative-HS}
\end{equation}
Here $E_T^{\rm geom}$ is the geometric edge operator of
Proposition~\ref{prop:geometric-edge-ledger}; it contains no part of
$H_{\rm ar}$. The regular component $R_T^{\rm edge+arch}$ is defined by
\eqref{eq:logonly-final-decomposition}, so its assumed relative bound must
also account for the entire $J_T^{\rm arch}$ block. If the geometric ledger
gives $L_T^{\rm geom}=E_T^{\rm geom}+R_T^{\rm geom}$ exactly, then
$R_T^{\rm edge+arch}=R_T^{\rm geom}+J_T^{\rm arch}$; the bound for
$R_T^{\rm geom}$ alone does not imply the bound assumed here. Put
$R_T^{\rm hor}:=R_T^{\rm hor,H}$. The regular operator is
\[
R_T^{\rm reg}:=R_T^{\rm edge+arch}+R_{f',T}^{\log}+R_T^{\rm hor}.
\]
Thus the exact four-block identity yields the conditional three-term equality
$A_F=A_T^{\rm true}+E_T^{\rm geom}+R_T^{\rm reg}$ on $\Pret$. The current
route has not proved the relative estimate for $R_T^{\rm edge+arch}$.
Higher stationary and Stirling terms
occur only in estimates of these exact operators.  Under
$\lambda_T\ge\lambda_*>0$, three separate unscaled relative
Hilbert--Schmidt estimates $o(N_T)$ imply the displayed scaled estimate by the
Hilbert--Schmidt triangle inequality.  No auxiliary scalar approximation or
frozen-profile defect is included: such terms would change the true operator
and are therefore retained inside $A_T^{\rm true}$.
In particular,
\begin{equation}
\left\|\Gfr_T^{-1/2}R_T^{\rm reg}\Gfr_T^{-1/2}\right\|_{S_2}^2
=o(\lambda_T^2N_T).
\label{eq:logonly-total-relative-HS}
\end{equation}
\end{hypothesis}
